\chapter{LINEAR ALGEBRA AND MULTILINEAR ALGEBRA.}

Linear algebra is at the same time one of the oldest branches of mathematics and one of the newest. On the one hand, problems are found in the origins of mathematics which are solved by a single multiplication or division, that is to say by the calculation of a value of a function \(f(x) = ax\) , or by the solution of an equation \(ax = b\): those are typical problems of linear algebra, and it is not possible to deal with them, nor even to state them correctly, without ``thinking linearly''.

On the other hand, not only these questions, but almost everything that touches on equations of the first degree, had for a long time been relegated to elementary teaching, when the modern development of the notions of field, ring, topological vector space, etc. . . ., came to free and give importance to the essential notions of linear algebra (for example duality); it is then that the essentially linear character of almost all modern algebra was realised, of which this ``linearisation'' is even one of the distinguishing traits, and that the place that is its due was given back to linear algebra. To give its history, from the point of view where we put ourselves, will thus be a task that is as important as it is difficult; and we will have to be content here with fairly summary indications.

From what comes earlier, it follows that linear algebra has without doubt been born in order to supply the needs of practising calculators; it is thus that we see the rule of three ({[}311{]}, v. I, pp.~150-155) and the rule of false position, more or less clearly stated, playing an important role in all the manuals of practical arithmetic, from the Rhind papyrus of the Egyptians up to those which are in use in our primary schools, passing through Aryabhata, the Arabs, Leonardo of Pisa and the innumerable ``books of calculation'' of the Middle Ages and of the Renaissance; but they have perhaps never been anything other than an extract, for the use of practicians, of more advanced scientific theories.

As for actual mathematicians, the nature of their research on linear algebra is a function of the general structure of their science. Ancient Greek mathematics, as it is set out in the Elements of Euclid, has developed two abstract theories of a linear character, on the one hand that of quantities ({[}107{]}, book V, cf.~p.~152), on the other hand that of whole numbers ({[}107{]}, book VII). With the Babylonians, we find methods that are much closer to our elementary algebra; they know how to solve, and in a very elegant way ({[}232{]}, pp.~181-183) systems of equations of the first degree. For a very long time, however, progress in linear algebra was mainly tied to that of algebraic calculation, and it is under this aspect, foreign to the present Note, that it is reasonable to consider it; to reduce, indeed, a linear system to an equation of the type \(ax = b\) , it is enough, if only one unknown is involved, to know the rules (already substantially stated by Diophantus) by which one can move terms from one member to another, and combine similar terms; and, if many unknowns are involved, it is enough to know as well how to eliminate them successively until only one remains. So the treatises on algebra, until the XVIIIth century, were thought to have done everything as far as the first degree is concerned, once they had laid out these rules; as for a system with as many unknowns as equations (they do not consider any others) where the first members would not be linearly independent forms, they content themselves invariably with noting in passing that that indicates a badly posed problem. In the treatises of the XIXth century, and even certain more recent works, this point of view has only been changed by the progress of notation, which allows the writing of systems of \(n\) equations in \(n\) unknowns, and by the introduction of determinants, which allow the statement of formulae for the explicit solution in the ``general case''; this progress, for which the credit would be due to Leibniz ({[}198 a{]}, v. I, p.~239) if he had published his ideas on the subject, is due principally to mathematicians of the XVIIIth and the beginning of the XIXth century.

But we must first consider several streams of ideas which, much more than the study of systems of linear equations, have contributed to the development of linear algebra in the sense in which we understand it. Inspired by the study of Apollonius, Fermat ({[}109{]}, v. I, pp.~91-110; French translation, v. III, pp.~85-101), conceives, even before Descartes ({[}85 a{]}, v. VI), the principles of analytic geometry, has the idea of the classification of plane curves by means of their degree (which, having become more and more familiar to all mathematicians, can be considered as definitely acquired towards the end of the XVIIth century), and states the fundamental principle of an equation of the first degree, in the plane, representing a straight line, and an equation of the second degree a conic: a principle from which he deduces straightaway some ``very pretty'' consequences concerning geometric loci. At the same time, he announces ({[}109{]}, v. I, pp.~184-188; French translation, v. III, pp.~159-163) the classification of problems as problems solved, problems which reduce to an equation in two unknowns, to an equation in three unknowns, etc.; and he adds: the first consist in determining a point, the second a line or plane locus, those following a surface; etc. (``\ldots such a problem does not seek only a point or a line, but a whole surface proper to the question; it is from that that are born the superficial loci, and so on for those following'', loc. cit. p.~186; there can be found already the embryo of geometry in n dimensions). This piece of writing, stating the principle of dimension in algebra and in geometric algebra, indicates a fusion of algebra and geometry, absolutely in line with modern ideas, but which, as has been seen already, took more than two centuries to penetrate minds.

At least these ideas ended up soon in the flowering of analytic geometry, which acquires its whole range in the XVIIIth century with Clairault, Euler, Cramer, Lagrange and many others. The linear character of the formulae for the transformation of co-ordinates in the plane and in space, which cannot fail to have been noticed already by Fermat, is put into prominence for example by Euler ({[}108 a{]}, (1), v. IX, chap.~II-III, and Append. chap.~IV), who builds on this the classification of planar lines, and that of surfaces, according to their degree (invariant precisely because of the linearity of these formulae); it is also he (loc. cit., chap.~XVIII) who introduces the word ``affinity'' to indicate the relation between curves that can be deduced from one another by a transformation \(x' = ax, y' = by\) (but without noticing anything of a geometric invariant in this definition which remains attached to a particular choice of axes). A bit later we see Lagrange ({[}191{]}, v. III, pp.~661-692) devoting a whole memoir, which will remain justly famous for a long time, to typically linear and multilinear problems of analytical geometry in three dimensions. It is at about the same time, in connection with the linear problem consisting of the search for a plane curve passing through given points, that the notion of determinant takes shape, first in a somewhat empirical manner, with Cramer {[}72{]} and Bézout {[}21{]}; developed afterwards by several authors, this notion and its essential properties are definitively brought out by Cauchy ({[}56 a{]}, (2), v. I, pp.~91-169) and Jacobi ({[}171{]}, v. III, pp.~355-392).

On the other hand, while mathematicians had some slight tendency to hold first degree equations in some disdain, the solution of differential equations constituted on the contrary a major problem; it was natural that, amongst these equations, one picked out at an early stage the linear equations, with constant coefficients or not, and that their study contributed to enhance the value of linearity and that which goes with it. That is just what is seen with d'Alembert {[}75 c{]}, Lagrange ({[}191{]}, v. I, pp.~471-668) and Euler ({[}108 a{]}, (1), v. XII); but only the first of these judges it useful to say that the general solution of a non homogeneous equation is the sum of a particular solution and of the general solution of the corresponding homogeneous equation; moreover, when these authors state that the general solution of the linear homogeneous equation of order n is a linear combination of n particular solutions, they do not add that these must be linearly independent, and make no effort to render explicit this latter notion; on this point as on so many others, it is only, it appears, the teaching of Cauchy at the École Polytechnique that will throw some light ({[}56 b{]}, pp.~573-574). But already Lagrange (loc. cit.) introduces also (solely for calculations, it is true, and without giving it a name) the adjoint equation \(L^{*}(y) = 0\) of a linear differential equation \(L(y) = 0\) , a typical example of duality by virtue of the relation

\[
\int z L (y) d x = \int L ^ {*} (z) y d x
\]

valid when y and z are zero at the ends of the interval of integration; more precisely, and 30 years before Gauss had explicitly defined the transpose of a linear substitution in 3 variables, we see there without doubt the first example of a ``functional operator'' \(L^{*}\) the transpose or ``adjoint'' of an operator L given by means of a bilinear function (here the integral \(\int yzdx\) ).

At the same time, and with Lagrange also ({[}191{]}, v. III, pp.~695-795), linear substitutions, in 2 or 3 variables at first, were on the way to conquer arithmetic. It is clear that the set of values of a function \(F(x, y)\) , when \(x\) and \(y\) take all integer values, does not change when an arbitrary linear substitution, with integer coefficients, of determinant 1, is performed on \(x, y\) ; it is on this fundamental observation that Lagrange bases the theory of the representation of numbers by forms, and that of the reduction of forms; and Gauss, by means of a step of which it has become difficult for us to appreciate all the boldness, gets out of it the notion of equivalence and that of a class of forms (cf.~p.~50); on this point, he recognises the necessity of some elementary principles relative to linear substitutions, and introduces in particular the notion of transpose or adjoint ({[}124 a{]}, v. I, p.~304). From that moment on, the arithmetic study and the algebraic study of quadratic forms in 2, 3 and later \(n\) variables, that of bilinear forms that are closely linked to them, and more recently the generalisation of these notions to an infinity of variables, was, up to our time, to form one of the most fruitful sources of progress for linear algebra (cf.~pp.~136 ff.).

But, a still more decisive progress perhaps, Gauss, in these same Disquisitiones, created (cf.~p.~51) the theory of finite abelian groups, which arises in four different ways, by means of the additive group of whole numbers modulo m (for m an integer), by means of the multiplicative group of numbers prime to m modulo m, by means of the group of classes of binary quadratic forms, and finally by means of the multiplicative group of mth roots of unity; and, as we have already noted, it is clearly as abelian groups, or to put it better as modules over Z, that Gauss treats all these groups, studies their structure, their isomorphism relations, etc.. In the module of ``complex integers'' \(a + bi\) , he is seen later studying an infinite module over Z, of which he has no doubt not missed seeing the isomorphism with the module of periods (discovered by him in the complex domain) of elliptic functions; in any case this idea appears already clearly with Jacobi, for example in the famous proof of the impossibility of a function with 3 periods and in his views on the problem of inversion of abelian integrals ({[}171{]}, v. II, pp.~25-50), to end up soon with the theorems of Kronecker ({[}186 a{]}, v. III \(_{1}\) , pp.~49-109).

Here, in streams of which we have sought to follow the thread and sometimes the meanderings, another comes to join in, which for a long time had remained underground. As it will be brought out elsewhere in more detail (see p.~130), ``pure'' geometry in the sense in which it had been understood in the previous century, that is to say essentially the projective geometry of the plane and of space without the use of co-ordinates, had been created in the XVIIth century by Desargues {[}84{]}, whose ideas, appreciated at their true worth by a Fermat, put into use by a Pascal, had subsequently fallen into disuse, eclipsed by the brilliant progress of analytical geometry; it is brought back to pride of place towards the end of the XVIIIth century, by Monge, then Poncelet and his emulators, Brianchon, Chasles, sometimes completely and voluntarily separated from analytical methods, sometimes (especially in Germany) mixed intimately with them. Yet projective transformations, from whatever point of view they are considered (synthetic or analytical), are none other of course than linear substitutions on the projective or ``barycentric'' co-ordinates; conics (in the XVIIth century), and later quadrics, of which the projective theory has for a long time been the principal object of study of this school, are none other than quadratic forms, of which we have already brought out above the close connection with linear algebra. To these notions is joined that of polarity: created, it also, by Desargues, the theory of poles and polars becomes in the hands of Monge and his successors, and soon under the name of principle of duality, a powerful instrument of transformation for geometric theorems; if one does not dare affirm that its links with adjoint differential equations were noticed only belatedly (they are indicated by Pincherle at the end of the century), what has none the less not been missed, Chasles bears testimony to it ({[}60 b{]}, p.~55), is to notice its kinship with the notion of reciprocal spherical triangle, introduced into spherical trigonometry by Viète ({[}319{]}, p.~418) and Snellius already in the XVIth century. But duality in projective geometry is only one aspect of the duality of vector spaces, taking account of the modifications that are imposed by passing from an affine space to a projective space (which is a quotient space of it, by the relation ``scalar multiplication'').

The XIXth century, more than any period in our history, was rich in mathematicians of the first rank; and it is difficult in a few pages, even limiting oneself to the most outstanding aspects, to describe all that the fusion of these movements of ideas produces in their hands. Between the purely synthetic methods on the one hand, a kind of Procrustean bed where their orthodox adepts torture themselves, and the analytic methods linked to a system of co-ordinates arbitrarily inflicted on space, one feels soon the need for a kind of geometrical calculus, dreamed of but not created by Leibniz, imperfectly outlined by Carnot {[}51{]}: it is first the addition of vectors that appears, implicit in Gauss in his geometrical representation of imaginary numbers and the application that he makes of it to elementary geometry (cf.~p.~162), developed by Bellavitis under the name of ``method of equipollence'', and that takes its definitive shape in Grassmann, Möbius, Hamilton; at the same time, under the name of ``barycentric calculus'', Möbius gives a version of it adapted to the needs of projective geometry ({[}223{]}, v. I).

At the same time, and amongst the same men, the passage takes place, so natural (as soon as this path has been entered) that we have seen it already announced by Fermat, from the plane and ``ordinary'' space to space of n dimensions; a passage that is even inevitable, since algebraic phenomena, which are interpreted for 2 or 3 variables as themselves in geometric terms, subsist without change for any number of variables; to impose therefore, in the use of geometric language, a limitation to 2 or 3 dimensions, would be for the modern mathematician a yoke as inconvenient as that which continually prevented the Greeks from extending the notion of number to ratios of incommensurable quantities. Also the language and ideas relating to space of n dimensions appear almost simultaneously on all sides, obscurely in Gauss, clearly among the mathematicians of the subsequent generation; and their more or less great measure of boldness in using it took less account of their mathematical leanings than of their philosophical views or even just practical considerations. In any case, Cayley and Grassmann, in about 1846, manipulate these concepts with the greatest of ease (and that, says Cayley differing from Grassmann ({[}58{]}, v. I, p.~321), ``without recourse to any metaphysical notion''); in Cayley, one remains constantly very near to the analytic interpretation and co-ordinates, whereas in Grassmann it is from the first, with the addition of vectors in space of n dimensions, the geometric aspect that takes the lead, to end up with developments of which we shall speak in a moment.

However, the impulsion that came from Gauss was pushing mathematicians, in two different ways, towards the study of algebras or hypercomplex systems. On the one hand, one could not avoid trying to extend the domain of real numbers other than by the introduction of the ``imaginary unit'' \(i = \sqrt{-1}\) , and to open up thus perhaps domains vaster and more fruitful than that of the complex numbers. Gauss convinced himself ({[}124 a{]}, v. II, p.~178) of the impossibility of such an extension, at least if one seeks to keep the principal properties of the complex numbers, that is to say in modern language, those that make it a commutative field; and, either because of his influence, or independently, his contemporaries seem to have shared this conviction, which will only be justified much later by Weierstrass ({[}329 a{]}, v. II, pp.~311-332) in a precise theorem. But, as soon as one interprets the multiplication of complex numbers as rotations in the plane, one is led, if intending to extend this idea to space, to envisage (since rotations in space form a non abelian group) non commutative multiplications; it is that which is one of the ideas that guides Hamilton \(^{1}\) in his discovery of quaternions {[}145 a{]}, the first example of a non commutative field. The singularity of this example (the only one, as Frobenius was to prove later, that can be constructed on the field of reals) restricted somewhat its scope, in spite or perhaps even because of the formation of a school of fanatical ``quaternionists'': a strange phenomenon, which repeats itself later around the work of Grassmann, then the popularisers who draw out of Hamilton and Grassmann what has been called the ``vector calculus''. The abandonment a little later of associativity, by Graves and Cayley who construct the ``Cayley numbers'', does not open up an interesting vein. But, after Sylvester had introduced matrices, and (without giving it a name) had clearly defined rank ({[}304{]}, v. I, pp.~145-151), it is again Cayley ({[}58{]}, v. II, pp.~475-496) who creates its calculus, not without noting (an essential fact, often lost sight of later on) that a matrix is only an abbreviated notation for a linear substitution, just as in fact Gauss denoted by \((a, b, c)\) the form \(aX^2 + 2bXY + cY^2\) . What is there is besides only one of the aspects, the most interesting for us no doubt, of the abundant production by Sylvester and Cayley on determinants and all that is linked with them, a production bristling with ingenious identities and impressive calculations.

It is also (among other things) an algebra over the reals that Grassmann discovers, the exterior algebra to which his name remains linked. His work, prior even to that of Hamilton ({[}134{]}, v I₁), created in a moral solitude which was almost complete, remained for a long time badly known, because no doubt of its originality, because also of the philosophical fog by which it begins to be shrouded, and which for example turns Möbius away at first. Furnished with concerns analogous to those of Hamilton but with a wider range (and which, as he realises soon, were those themselves of Leibniz), Grassmann constructs a vast algebraico-geometric edifice, resting on a geometric or ``intrinsic'' conception (already just about axiomatised) of the vector space of n dimensions; amongst the most elementary of the results which he ends up with, let us quote for example the definition of linear independence of vectors, that of dimension and the fundamental relation dim V + dim W = dim (V + W) + dim (V ∩ W) (loc. cit., p.~209; cf.~{[}134{]}, v. I₂, p.~21). But it is above all the exterior multiplication, then the interior multiplication, of multivectors, that provide him with the tools by means of which he deals easily with, first the problems of linear algebra proper, then those which are concerned with the Euclidean structure, that is to say with the orthogonality of vectors (where he finds the equivalent of the duality which he is missing).

The other path opened up by Gauss, in the study of hypercomplex systems, is that which starts with the complex integers \(a + bi\) ; from these, one carries on quite naturally to more general algebras or hypercomplex systems, over the ring Z of integers or over the field Q of rationals, and first of all to those foreseen already by Gauss which are generated by the roots of unity, then to the fields of algebraic numbers and to modules of algebraic integers: those are the principal objects of the work of Kummer, the study of the former is tackled by Dirichlet, Hermite, Kronecker, Dedekind. Here, contrary to what happens with algebras over the reals, it is not necessary to give up any of the characteristic properties of commutative fields, and it is to these that, during the whole of the XIXth century, work is limited. But the linear properties, and for example the search for a basis for the integers of the field (indispensable for a general definition of the discriminant) play in many respects an essential role; and, in Dedekind in any case, the methods are destined to become typically ``hypercomplex''; besides Dedekind himself, without addressing in a general way the problem of algebras, is aware of this aspect of his work, and of that which ties it for example to the results of Weierstrass on hypercomplex systems over the reals ({[}79{]}, in particular vol.~II, pp.~1-19). At the same time the determination of the structure of the multiplicative group of units in a field of algebraic numbers, carried out by Dirichlet in his famous notes ({[}92{]}, v. I, pp.~619-644) and almost at the same time by Hermite ({[}159{]}, v. I, pp.~159-160), was eminently suitable to clarify the ideas about modules over Z, their systems of generators, and, when they have them, their bases. Then the notion of ideal, that Dedekind defines in algebraic number fields (as a module over the ring of integers of the field) ({[}79{]}, v. III, p.~251), whereas Kronecker introduces in polynomial rings (under the name of ``system of modules'') an equivalent notion ({[}186 a{]}, v. II, pp.~334-342), gives the first examples of modules over rings more general than Z; and in the same authors, then in Hilbert, is released little by little in particular cases the notion of a group with operators, and the possibility of always constructing, starting with such a group, a module over a suitably defined ring.

At the same time, the arithmetico-algebraic study of quadratic and bilinear forms, and their ``reduction'' (or, what amounts to the same, that of matrices and their ``invariants'') leads to the discovery of general principles concerning the solution of systems of linear equations, principles that, for lack of the notion of rank, had escaped Jacobi. \(^{2}\) The problem of the solution in integers of systems of linear equations with integral coefficients is tackled and solved, at first in a special case by Hermite, then in all its generality by H. J. Smith ({[}287{]}, v. I, pp.~367-409); the results of this latter are found again, only in 1878, by Frobenius, within the framework of a vast programme of research instituted by Kronecker, and in which Weierstrass also participates; it is incidentally Kronecker, in the course of his work, who gives the definitive form to the theorems on linear systems with real (or complex) coefficients, which is also elucidated in an obscure manual, with the scrupulous care which characterises him, by the famous author of Alice in Wonderland; as for Kronecker, he disdains publishing his results, and leaves it to his colleagues and disciples; even the word rank is only introduced by Frobenius. It is also in their teaching at the university of Berlin that Kronecker {[}186 b{]} and Weierstrass introduce the ``axiomatic'' definition of determinants (as a multilinear alternating function of n vectors in space of n dimensions, normalised so as to take the value 1 for the identity matrix), a definition equivalent to that which can be deduced from the calculations of Grassmann; it is in his lectures that Kronecker still, without feeling the need to give it a name, and in a still not intrinsic form, introduces the tensor product of spaces and the ``Kroneckerien'' product of matrices (a linear substitution induced, on a tensor product, by given linear substitutions applied to the factors).

These researches should not be separated either from invariant theory, created by Cayley, Hermite and Sylvester (the ``invariant trinity'' of whom later Hermite speaks in his letters) and which, from a modern point of view, is above all a representation theory of linear groups. There emerges, as an algebraic equivalent of the duality of projective geometry, the distinction between series of cogredient and contragredient variables, that is to say vectors in a space and vectors in the dual space; and after attention had been drawn first to forms of low degree, then of arbitrary degree, in 2 or 3 variables, little time is lost in examining bilinear, then multilinear forms, in several series of ``cogredient'' or ``contragredient'' variables, which is equivalent to the introduction of tensors; this becomes conscious and is popularised when, under the inspiration of invariant theory, Ricci and Levi-Civita, in 1900, introduce the ``tensor calculus'' into differential geometry {[}257{]}, which later was very much in fashion following its use by ``relativist'' physicists. It is still the progressive interpenetration of invariant theory, of differential geometry, and of the theory of partial differential equations (especially the problem of Pfaff so called, and its generalisations) which brings little by little geometers to consider, first of all bilinear alternating forms of differentials, in particular the ``bilinear covariant'' of a form of degree 1 (introduced in 1870 by Lipschitz, then studied by Frobenius), to finish with the creation, by E. Cartan ({[}52 a{]}, v. II₁, pp.~303-396) and Poincaré ({[}251 b{]}, v. III, chap.~XXII) of the calculus of exterior differential forms. Poincaré introduces these, in the context of the formation of his integral invariants, as being expressions which appear in multiple integrals, whereas Cartan, guided no doubt by his researches on algebras, introduces them in a less formal way, but not without observing also that the algebraic part of their calculus is identical to the exterior multiplication of Grassmann (hence the name he adopts), thereby putting definitively into its true place the latter's work. The translation into the notation of tensor calculus, of exterior differential forms, shows besides immediately their link with asymmetric tensors, which will, as soon as one returns to a purely algebraic point of view, make it to be seen that they are to alternating multilinear forms what covariant tensors are to arbitrary multilinear forms; this aspect becomes clearer still with the modern theory of representations of linear groups; and by that is recognised, for example, the substantial identity between the definition of determinants given by Weierstrass and Kronecker, and that which results from the calculus of Grassmann.

We reach thus the modern era, where the axiomatic method and the notion of structure (felt at first, defined only recently), allow us to separate concepts that up till then had been inextricably mixed, to formulate that which was vague or subconscious, to prove in the generality which is appropriate the theorems which had only been known in particular cases. Peano, one of the creators of the axiomatic method, and one of the first mathematicians also to appreciate at its true worth the work of Grassmann, gives already in 1888 ({[}246 b{]}, chap.~IX) the axiomatic definition of vector spaces (whether of finite dimension or not) over the field of reals, and, with a fully modern notation, linear maps from such a space to another; a little later, Pincherle seeks to develop applications of linear algebra, thus conceived, to the theory of functions, in a direction, it is true, which has hardly shown itself to be fruitful; at least his point of view will allow him to recognise in the ``Lagrange adjoint'' a particular case of the transposition of linear maps: which appears soon, still more clearly, and for partial differential equations as well as for differential equations, in the course of the memorable work of Hilbert and his school on Hilbert space and its applications to analysis. It is in connection with these latter researches that Toeplitz {[}309 a{]}, also introducing (but by means of co-ordinates) the most general vector space over the reals, makes the fundamental observation that the theory of determinants is useless for the proof of the main theorems of linear algebra, which allows the extension of these without difficulty to infinite dimensional spaces; and he indicates also that the linear algebra thus understood is applicable of course to any commutative base field.

On the other hand, with the introduction by Banach, in 1922, of the spaces which carry his name, \(^{3}\) one meets, it is true in a problem which is topological as much as algebraic, spaces which are not isomorphic to their dual (see pp.~214 ff.). Already between a vector space and its dual, there is no ``canonical'' isomorphism, that is to say determined by the structure, which had been reflected for a long time by the distinction between cogredient and contragredient. It seems nonetheless that there is no doubt that the distinction between a space and its dual was only definitively established following the work of Banach and his school; it is this work too which raised an interest in the notion of codimension. As for the duality or ``orthogonality'' between vector subspaces of a space and those of its dual, the way in which it is formulated today demonstrates an analogy which is not just coincidental with the modern formulation of the main theorem of Galois theory, and with the duality of Pontrjagin so called in locally compact abelian groups; this goes back to Weber, who, in the course of arithmetic researches, stated the bases of it in 1866 for finite groups {[}327 d{]}; in Galois theory, the ``duality'' between subgroups and subfields takes shape with Dedekind and Hilbert; and the orthogonality between vector subspaces is visibly derived, first of all from the duality between linear varieties in projective geometry, and also from the notion and properties of completely orthogonal varieties in a Euclidean space or a Hilbert space (whence its name). All these threads gather together in the contemporary era, in the hands of algebraists such as E. Noether, Artin and Hasse, of topologists such as Pontrjagin and Whitney, not without mutual influence being exerted between one group and the other.

At the same time a critical examination is being undertaken, which is meant to eliminate, at each point, the hypotheses which are not truly indispensable, and especially those which would prevent certain applications. Thus is realised the possibility of substituting rings for fields in the notion of vector space, and, creating the general notion of module, treating at one and the same time these spaces, abelian groups, particular modules already studied by Kronecker, Weierstrass, Dedekind, Steinitz, and even groups with operators, and for example to apply to them the Jordan-Hölder theorem; at the same time the transition to non commutativity is made, by means of the distinction between right modules and left modules, to which the modern development of the theory of algebras by the American school (Wedderburn, Dickson) and especially the German school (E. Noether, Artin) was leading.

Finally there has emerged, quite recently, the last of the tendencies of which we must speak here: contained in embryo in the theorem of Dedekind ({[}79{]}, vol.~III, p.~29) according to which arbitrary automorphisms of a commutative field are always linearly independent, the linearisation of Galois theory is carried out with Artin {[}7 a{]}, then it is soon generalised to arbitrary extensions of commutative fields, and even to extensions of non commutative fields ({[}172 d{]}, chap.~VII).

\chapter{POLYNOMIALS AND COMMUTATIVE FIELDS.}

The theory of commutative fields --- and the theory of polynomials, which is closely linked to it --- derive directly from what constituted, until the middle of the XIXth century the principal object of classical Algebra: the solution of algebraic equations, and geometric construction problems which are equivalent to that.

As soon as the solution of an algebraic equation of degree \textgreater{} 1 is sought, one finds oneself up against calculation difficulties of a totally new kind, the determination of an unknown being impossible to accomplish by means of ``rational'' calculations starting from given quantities. This difficulty must have been seen at a very early stage: and it must be counted as one of the most important contributions of the Babylonians that they knew how to reduce the solution of quadratic and biquadratic equations to a single new algebraic operation, the extraction of square roots (as is proved by the numerous numerical equations that are found solved in this way in the texts that have reached us ({[}232{]}, pp.~183-193)). As far as formal calculation is concerned, Antiquity will never reach further than this point in the problem of solving algebraic equations; the Greeks of the classical era limit themselves indeed to finding again the Babylonian formulae in geometrical terms, and their use in an algebraic form is not attested to before Hero (100 after J. C.) and Diophantus.

It is in a totally different direction that the Greeks make decisive progress. We have poor information on the way in which the Babylonians conceived of and calculated the square roots of non square integers; \(^{1}\) in the rare texts which have reached us on this question, they seem to be satisfied with fairly crude approximation methods ({[}232{]}, pp.~33-38). The Pythagorean school which had determined with rigour the concept of commensurable quantities, and attached to it a quasi religious character, could not hold this point of view; and it is possible that it was the failure of repeated attempts to express \(\sqrt{2}\) as a rational that led them finally to prove that this number is irrational. \(^{2}\)

We will say elsewhere (cf.~p.~148) how this discovery, which marks a major turning point in the history of Mathematics, reacted profoundly on the concept of ``number'' among the Greeks, and led them to create an algebra of exclusively geometric character, so as to find in it a way of representing (or perhaps a proof of ``existence'' for) the incommensurable ratios, that they refused to consider as numbers. Most often, an algebraic problem is reduced by them to the intersection of two auxiliary plane curves, suitably chosen, or to several successive determinations of such intersections. Late traditions with little authority make the introduction of a first classification of these constructions, destined for a long and brilliant career, go back to Plato: for reasons that are more philosophical than mathematical, it appears, he would have put aside the constructions ``by ruler and compass'' so called, that is to say those where there are no auxiliary curves except for straight lines and circles. \(^{3}\) In any case, Euclid, in his Elements {[}107{]}, limits himself exclusively to treating problems soluble in this way (without however characterising them by a particular name); a circumstance which doubtless would contribute not a little to fastening the attention of mathematicians of the following centuries on these problems. But we know today \(^{4}\) that algebraic equations that can be solved ``by ruler and compass'' are of a very special type; in particular, an irreducible equation of the 3rd degree (over the field of rationals) can not be solved in this way, and the Greeks had met very early on problems of the 3rd degree which have remained famous, such as the duplication of the cube (solution of \(x^{3}=2\) ) and the trisection of an angle; the quadrature of the circle put them on the other hand in presence of a transcendental problem. We see them, in order to solve these problems, introducing numerous algebraic (conics, cissoid of Diocles, conchoid of Nicomedes) or transcendental (quadratic of Dinostratus, spiral of Archimedes) curves; which could not have failed to lead them to an autonomous study of these various curves, thus preparing the way for future developments in analytic Geometry, in algebraic Geometry and in infinitesimal Calculus. But these methods do not help to make any progress in the solution of algebraic equations, \(^{5}\) and the only work of Antiquity which made a worthy contribution to this question and exerted a lasting influence on the algebraists of the Middle Ages and the Renaissance is Book X of the Elements of Euclid {[}107{]}; in this Book (to which some historians would make the principal results of Theetetus go back), he considers the ex-

pressions obtained by the combination of several radicals, such as \(\sqrt{\sqrt{a}} \pm \sqrt{b}\) (a and b rational), gives conditions under which these expressions are irrational, classes them into numerous categories (which he shows are distinct), and studies the algebraic relations between these various irrationals, such as those which we would write today as

\[
\sqrt {\sqrt {p} + \sqrt {q}} = \sqrt {1 / 2 (\sqrt {p} + \sqrt {p - q})} + \sqrt {1 / 2 (\sqrt {p} - \sqrt {p - q})};
\]

the whole expressed in the usual geometric language of the Elements, which makes the exposition particularly heavy and awkward.

After the decline of classical Greek mathematics, the concepts concerning algebraic equations are modified. There is no doubt that, during the whole of the classical period, the Greeks were in possession of methods of indefinite approximation for square roots, about which we are unhappily badly informed. \(^{6}\) With the Hindus, then the Arabs and their western emulators in the Middle Ages, the extraction of roots of all orders becomes an operation that tends to be considered as fundamental on the same basis as the rational operations of algebra, and to be denoted, like these latter, by symbols which are more and more manageable in calculations. \(^{7}\) The theory of the equation of the second degree, which is perfected during the whole of the Middle Ages (number of roots, negative roots, cases of impossibility, double roots), and that of biquadratic equations, give models for formulae for the solution of equations ``by radicals'', from which algebraists are going to try for centuries to copy analogous formulae for the solution of equations of higher degree, and in the first instance for equations of the 3rd degree. Leonardo of Pisa, the principal introducer of Arab science to the West in the XIIIth century, recognises in any case that the irrationals classified by Euclid in his Xth book can not be of use to this end (a new proof of impossibility, in a theory that contains so many), and we see him already trying to do similar calculations with cube roots, obtaining relations such as

\[
\sqrt [ 3 ]{16} + \sqrt [ 3 ]{54} = \sqrt [ 3 ]{250},
\]

analogous to Euclid's formulae for square roots (and of which besides previous examples are found among the Arabs). But three centuries of unfruitful efforts will still go by before Scipio del Ferro, at the beginning of the XVIth century, finally succeeds, for the equation \(x^{3} + ax = b\) , in obtaining the solution formula \(^{8}\)

\[
x = \sqrt [ 3 ]{\frac {b}{2} + \sqrt {\left(\frac {b}{2}\right) ^ {2} + \left(\frac {a}{3}\right) ^ {3}}} + \sqrt [ 3 ]{\frac {b}{2} - \sqrt {\left(\frac {b}{2}\right) ^ {2} + \left(\frac {a}{3}\right) ^ {3}}}.\tag{5.1}
\]

We cannot describe here the picturesque side of this sensational discovery --- the quarrels that it provoked between Tartaglia on the one hand, Cardan and his school on the other -- nor the figures, often appealing, of the scholars who were the protagonists. But we must note the decisive progress that ensued for the theory of equations, in the hands of Cardan and his pupils. Cardan, who had less repugnance than most of his contemporaries at using negative numbers, observes thus that equations of the third degree may have three roots, and the biquadratic equations four ({[}50{]}, v. IV, p.~259), and he remarks that the sum of the three roots of \(x^3 + bx = ax^2 + c\) (an equation in which he knows how to eliminate the term in \(x^2\) ) is always equal to \(a\) (ibid.). No doubt guided by this relation and the intuition of his general character, he has the initial idea of the notion of multiplicity of a root; above all he is emboldened (not without oratorical preparation) to calculate formally with expressions containing square roots of negative numbers. It is likely that he was brought to this by the fact that such expressions arise naturally in the use of the formula (1) when \(\left(\frac{b}{2}\right)^{2}+\left(\frac{a}{3}\right)^{3}<0\) (a case --- called ``irreducible'' --- where Cardan had recognised the existence of three real roots); it is what appears in any case with clarity in his disciple R. Bombelli, who, in his Algebra ({[}28 a{]}, p.~293) proves the relation

\[
\sqrt [ 3 ]{2 + \sqrt {-121}} = 2 + \sqrt {-1}
\]

and takes care to give explicitly the rules for the calculation of complex numbers in a form which is very close to modern expositions. \(^{9}\) Finally, already in 1545, another pupil of Cardan, L. Ferrari, succeeds in solving the general equation of the 4th degree, with the help of an auxiliary equation of the 3rd degree. \(^{10}\)

After such a rapid advance, the following period, until the middle of the XVIIIth century, does little more than develop the new ideas introduced by the Italian school. Thanks to the essential progress that it brings to algebraic notation, Viète can express in a general manner the relations between the coefficients and roots of an algebraic equation, at least when the roots are all positive \(^{11}\) ({[}319{]}, p.~158). Bolder, A. Girard {[}129{]} does not hesitate to affirm (of course without proof) that an equation of degree n has exactly n roots, on condition that ``impossible roots'' are counted, each with its degree of multiplicity, and that these roots satisfy the relations given by Viète; he obtains also, for the first time, an expression for the sums of powers of roots up to exponent 4.

But the spirit of the XVIIth century is turned towards other directions and it is only as an indirect consequence that Algebra profits somewhat from the new discoveries of analytic Geometry and the infinitesimal Calculus. So, to the method of Descartes for obtaining the tangents to algebraic curves (cf.~p.~178) is linked the criterion for the multiplicity of a root of an algebraic equation, which is stated by his disciple Hudde ({[}85 b{]}, v. I, pp.~433 and 507-509). It is no doubt also to the influence of Descartes that one must go back for the distinction between algebraic functions and transcendental functions, parallel to that which he introduces in his Géométrie between ``geometric'' and ``mechanical'' curves (cf.~pp.~176 and 193.) In any case, this distinction is perfectly clear in J. Gregory who, in 1667, tries even to prove that the area of a circular sector can not be an algebraic function of the chord and radius ({[}136 a{]}; cf.~{[}155{]}). The expression ``transcendental'' is from Leibniz, who does not cease to be interested in these questions of classification for the whole of his career, and who, in about 1682, discovers a simple proof of the result sought by Gregory, by proving that sin x can not be an algebraic function of x ({[}198 a{]}, v. V, pp.~97-98). \(^{12}\) With his friend Tschirnhaus, he is besides one of the only mathematicians of his time who is still interested in the problem of the solution ``by radicals'' of algebraic equations. In his early stages, we see him studying the ``irreducible case'' of the equation of the 3rd degree, and convincing himself (in fact without sufficient proof) that it is impossible in this case to get rid of imaginary quantities from the formulae of the solution ({[}198 d{]}, pp.~547-564). At about the same time, he attacks also without success the solution by radicals of the equation of the 5th degree; and when later, Tschirnhaus pretends to have solved the problem by making all the terms of the equation disappear except for the two extremes by means of a transformation of the form \(y = P(x)\) , where P is a polynomial of the 4th degree, suitably chosen, Leibniz realises straightaway that the equations which determine the coefficients of \(P(x)\) are of degree \textgreater{} 5, and judges that the method is destined for failure ({[}198 d{]}, pp.~402-403).

It appears that it is the needs of the new Analysis that had little by little revived interest in algebra. The integration of rational functions, accomplished by Leibniz and Johann Bernoulli, and the question of imaginary logarithms which is closely linked to it, are the occasion for deepening the calculus of imaginary numbers, and for answering the question of the decomposition of a polynomial into factors of the first degree (``fundamental theorem of algebra''). \(^{13}\) From the beginning of the XVIIIth century, the solution of the binomial equation \(x^{n}-1=0\) is reduced by Cotes and de Moivre to the division of the circle into n equal parts; in order to obtain the expression of its roots ``by radicals'', it is sufficient therefore to know how to do it when n is an odd prime, and de Moivre remarks that the substitution \(y = x + \frac{1}{x}\) reduces the problem then to the solution ``by radicals'' of an equation of degree \((n - 1)/2\) . As far as the ``fundamental theorem'' is concerned, after the repeated setbacks in the general solution ``by radicals'' (including several attempts by Euler ({[}108 a{]}, (1), v. VI, pp.~1-19 and 170-196)), proofs begin to be sought a priori, not using explicit formulae for the solution. Without going into the detail of the proposed methods (which would end up with the proofs of Lagrange and Gauss; cf.~p.~91 and 161), it is fitting to note here the point of view from which the problem is seen in the middle of the XVIIIth century: it is assumed (without any justification other than a vague feeling of generality, coming no doubt, as with A. Girard, from the existence of relations between the coefficients and roots) that an equation of degree n always has n ``ideal'' roots, with which one can calculate as with numbers, without knowing if they are numbers (real or complex); and what it is necessary to prove is (by using if necessary calculations on ideal roots) that at least one of these roots is an ordinary complex number. \(^{14}\) In this defective form, it is recognised that here is the first embryo of the general idea of a ``formal adjunction'' which, in spite of the objections of Gauss ({[}124 a{]}, v. III, p.~1), would become the basis of the modern theory of commutative fields.

With the fundamental memoirs of Lagrange ({[}191{]}, v. III, pp.~205-421) and of Vandermonde {[}315{]}, the year 1770 sees a new and decisive period in the history of the theory of algebraic equations opening up. To the empiricism of the more or less fortunate attempts at solution formulae, which until then had reigned without dispute, a systematic analysis of the problems as stated and of the methods susceptible to solve them is going to follow, an analysis which in sixty years will lead to the definitive results of Galois. Lagrange and Vandermonde both start from the ambiguity which multiple determinations of radicals introduce in the formulae for the solution of equations of degree \(\leq 4\) ; this fact had attracted the attention of Euler ({[}108 a{]}, (1), v. VI, pp.~1-19) who had showed among other things how, in the formula of del Ferro, one must associate the determinations of the radicals which appear there so as to obtain 3 roots and not 9. \(^{15}\) Lagrange remarks that each of the cubic radicals in the formula of del Ferro can be written in the form \(\frac{1}{3}(x_1 + \omega x_2 + \omega^2 x_3)\) , where \(\omega\) is a cube root of one, \(x_1, x_2, x_3\) are the three roots of the equation in question, taken in a certain order, and he makes the fundamental observation that the function \((x_1 + \omega x_2 + \omega^2 x_3)\) of the three roots can only take two distinct values for all permutations of the three roots, which explains a priori the success of the methods for the solution of this equation. A similar analysis of the solution methods for the equation of the 4th degree leads him to the function \(x_{1}x_{2} + x_{3}x_{4}\) of the four roots, which only takes three distinct values for all permutations of the roots, and is as a consequence a root of an equation of the third degree with coefficients that are rational functions of those of the given equation; \(^{16}\) these facts constitute, says Lagrange, ``the true principles, and, so to say, the metaphysic \(^{17}\) of the solution of equations of the 3rd and of the 4th degree'' ({[}191{]}, v. III, p.~357). Relying on these examples, he proposes studying in general, for an equation of degree n, the number \(\nu\) of values \(^{18}\) that a rational function V of the n roots can take when they are arbitrarily permuted; in reality he inaugurates thus (with this terminology still tightly adapted to the theory of equations) the theory of groups and that of fields, of which he obtains already several fundamental results by use of the same principles as those that are used today. For example, he shows that the number \(\nu\) is a divisor of n! by means of the argument that is used today for proving that the order of a subgroup of a finite group divides the order of the group. More remarkable still is the theorem where he shows that, if \(V_{1}\) and \(V_{2}\) are two rational functions of the roots such that \(V_{1}\) and \(V_{2}\) are invariant under the same permutations, then each of them is a rational function of the other and the coefficients of the equation (a particular case of the theorem of Galois characterising a subextension of a Galois extension as the field of invariants of its Galois group): ``This problem'', he says, ``appears to me one of the most important in the theory of equations, and the general solution that we will give for it will be the means of throwing new light on this part of Algebra'' ({[}191{]}, v. III, p.~374).

All these researches are naturally, in the mind of Lagrange, preliminaries to the analysis of possible methods for the solution of algebraic equations by a successive reduction to equations of lesser degree, such a method being linked, as he shows, to the formation of rational functions of the roots taking less than n values under permutations of the roots. Guided no doubt by his results on the equation of the 3rd degree, he introduces in general the ``Lagrange resolvents'' \(y_{k} = \sum_{h=1}^{n} \omega_{k}^{h} x_{h}\) , where \(\omega_{k}\) is an nth root of unity ( \(1 \leq k \leq n\) ), shows clearly how the knowledge of these n numbers results in that of the roots \(x_{k}\) , and looks in general for the degree of the equation which is satisfied by the \(y_{k}\) ; he shows for example that, if n is prime, the \(y_{k}\) are roots of an equation of degree n-1, whose coefficients are rational functions of a root of an equation of degree \((n-2)!\) with coefficients that are expressed rationally in terms of the coefficients of the given equation. ``Here are, if I am not mistaken'', he concludes, ``the real principles for the solution of equations, and the most correct analysis for getting there; everything reduces, as can be seen, to a kind of calculation of combinations, by which one finds a priori the results which should be expected'' ({[}191{]}, v. III, p.~403).

As for the memoir of Vandermonde, independent of that of Lagrange, it meets this latter at many points, particularly so far as the idea of seeking rational functions of the roots taking as few distinct values as possible under permutations of the roots is concerned, \(^{19}\) and the study of the ``Lagrange resolvents'' that he also introduces for this purpose. His work is far from having the clarity and generality of that of Lagrange; on one point however it goes clearly further, in applying the same ideas to the equation for the division of the circle \(x^{n}-1=0\) for n an odd prime. While Lagrange is content to remind readers that this equation reduces to an equation of degree \(m=(n-1)/2\) , with rational coefficients, without seeking to solve it when \(n\geq11\) , Vandermonde states that the mth powers of the Lagrange resolvents of this equation are rational, because of the relations between the various roots of \(x^{n}-1=0\) ; but he limits himself to checking the soundness of this assertion in the case n=11, without justifying it in a general way.

It is only 30 years later that the result announced by Vandermonde was fully proved by C. F. Gauss. \(^{20}\) His decisive results on the equation \(x^n - 1 = 0\) ( \(n\) an odd prime) fit into the general programme of his memorable research on arithmetic ({[}124 a{]}, v. I, pp.~413 ff.), and illustrate in a very special way his mastery in the handling of what we would now call the theory of cyclic groups. After having proved that the polynomial \(\Phi_n(x) = (x^n - 1)/(x - 1)\) is irreducible for \(n\) an odd prime, \(^{21}\) he has the idea of writing its \(n - 1\) roots in the form \(\zeta^{g^k} = \zeta_k (0 \leq k \leq n - 2)\) , where \(g\) is a primitive root of the congruence \(z^{n-1} \equiv 1 \pmod{n}\) ) (which amounts, in modern language, to bringing into the open the fact that the group \(\Gamma\) of the equation \(\Phi_n(x) = 0\) is cyclic). To every divisor \(e\) of \(n-1\) , he puts in correspondence the \(f = (n-1)/e\) ``periods'' \(\eta_\nu = \xi_\nu + \xi_{\nu+f} + \xi_{\nu+2f} + \cdots + \xi_{\nu+(e-1)f} (1 \leq \nu \leq f)\) and shows substantially that the linear combinations with rational coefficients of the \(\eta_{\nu}\) make up a field, generated by any one of the \(f\) periods \(\eta_{\nu}\) and of degree \(f\) over the field of rational numbers (this field corresponds naturally to the subgroup of \(\Gamma\) of order \(e\) ). We cannot go into the details of his analysis here, and of the important arithmetical consequences that follow from it; let us notice only that it gives him in particular the famous theorem on the possibility of constructing ``by ruler and compass'' the polygons having a number of sides equal to a prime number of the form \(2^{2^k} + 1\) . As for the solution by radicals of the equation \(\Phi_n(x) = 0\) , it follows easily from the theory of periods, applied to the \(f\) th power of a Lagrange resolvant \(\sum_{\nu=0}^{f-1} \omega^\nu \eta_\nu\) (where \(\omega^f = 1\) ).

It is directly to Lagrange that the research of his compatriot Ruffini {[}265{]} is linked, contemporaneous with the Disquisitiones; taking up again the question at the point at which Lagrange had left it, he proposes as goal the proof of the impossibility of the solution ``by radicals'' of the ``general'' equation \(^{24}\) of the 5th degree. The proof of Ruffini, prolix and obscure, remains incomplete although reworked several times; but it is already very close to the proof (correct in principle) that Abel would obtain later. \(^{25}\) His main interest resides especially in the introduction of calculations with substitutions and the first notions of the theory of groups, that Ruffini develops to show that there does not exist a function of the 5 roots of the equation taking more than 2 and less than 5 values when the roots are arbitrarily permuted.

We have already said (see p.~51) how this first sketch of the theory of groups of permutations was developed and systematised by Cauchy a few years later. But, if the notions necessary for the development of the ideas of Lagrange are thus clarified little by little as far as substitutions are concerned, it was still necessary to state in as clear a way the first principles of the theory of fields. That was what was missing for Ruffini, and it is that which Abel and Galois are going to do, in the last phase of the problem of the solution of algebraic equations.

During the whole of his short life, Abel does not fail to be concerned with this problem. While almost still a child, he believed that he had obtained a formula for the solution by radicals of the general equation of the 5th degree. Having realised later his mistake, he did not rest until he had succeeded in proving that such a formula did not exist ({[}1{]}, v. I, p.~66). But he does not stop there. While his emulator Jacobi develops the theory of elliptic functions as an analyst, it is the algebraic point of view that dominates the work of Abel on this question, centred on the theory of the equations of division of elliptic functions ({[}1{]}, v. I, pp.~265, 377 and passim). He obtains thus new types of equation soluble by radicals by a method copied from that of Gauss for the solution of equations of division of the circle ({[}1{]}, v. I, pp.~310 and 358); \(^{26}\) a result from which he rises to the conception of ``abelian'' equations, for which he proves, in a famous memoir, the solubility by radicals ({[}1{]}, v. I, p.~478); it is on this occasion that he defines in a precise way the notion of a polynomial irreducible over a given field (the field generated by the coefficients of the equation that he is studying). \(^{27}\) Finally death laid him low in 1829, just as he attacks the general problem of the characterisation of all the equations soluble by radicals, and having just communicated to Crelle and Legendre results already very near to those of Galois ({[}1{]}, v. II, pp.~219-243, 269-270 and 279).

It was to him that was reserved, three years later, the crowning of the structure {[}123{]}. Like Abel, but in a still clearer way, he starts by defining (up to terminology) the membership by a number of a field generated by given numbers, the notion of adjunction, and irreducible polynomials over a given field. Given an equation \(F(x) = 0\) , without multiple roots, with coefficients in a given field \(K\) , he shows successively that ``one can always construct a function \(V\) of the roots, such that none of the values that are obtained by permuting in this function the roots in all possible ways, are equal'', that this function ``enjoys the property that all the roots of the proposed equation can be expressed rationally as a function of \(V\)'', and that, \(V, V', V'', \ldots\) being the roots of the irreducible equation which \(V\) satisfies, ``if \(a = f(V)\) is a root of the proposed equation, \(f(V')\) as well will be a root of the proposed equation'' ({[}123{]}, pp.~47-51); in modern language, he thus proves that \(V\) , as well as any one of its conjugates over \(K\) , generates the field \(N\) of the roots of \(F\) . He then defines the group \(\Gamma\) of \(F\) as the set of permutations of the roots \(x_i\) that are obtained by substituting for \(V\) , in the rational expression for each of the \(x_i\) as a function of \(V\) , any one of the conjugates of \(V\) ; and he obtains straightaway the fundamental characterisation of the elements of \(K\) by the property of being invariant under every permutation of \(\Gamma\) ({[}123{]}, p.~51). He proves next that, if \(N\) contains the field of roots \(L\) of another polynomial, the group of \(N\) over \(L\) is a normal subgroup of \(\Gamma\) (a notion that he introduces at this point) ({[}123{]}, p.~175). From there he deduces finally the criterion of solubility of an equation by radicals, by means of an argument of which this is the essential part: assuming that the base field \(K\) contains all the roots of unity, there must exist by hypothesis an increasing sequence \((K_i)_{0 \leq i \leq m}\) of fields intermediate between \(K\) and \(N\) , with \(K_0 = K, K_m = N\) , the field \(K_{i+1}\) being obtained by adjoining to \(K_i\) all the roots of a binomial equation \(x^{n_i} - a_i = 0\) (with \(a_i \in K_i\) ). There exists therefore in \(\Gamma\) a decreasing sequence \((\Gamma_i)\) of subgroups such that \(\Gamma_0 = \Gamma, \Gamma_m = \varepsilon\) (neutral element), \(\Gamma_{i+1}\) being normal in \(\Gamma_i\) and the quotient group \(\Gamma_i / \Gamma_{i+1}\) being cyclic (a case in which one says that the group \(\Gamma\) is soluble). Conversely, if this is so, the use of a Lagrange resolvant shows that \(K_{i+1}\) is obtained by adjoining to \(K_i\) all the roots of a binomial equation, and hence the equation \(F(x) = 0\) is soluble by radicals. \(^{28}\) The impossibility of the solution by radicals of the ``general'' equation of degree \(n > 4\) is then a consequence of the fact that the group \(\Gamma\) of this equation, isomorphic to the symmetric group \(S_n\) is not soluble.

Starting with the middle of the XIXth century, the algebraists, as we have already remarked (cf.~p.~53), enlarge considerably the field of their investigations, until then almost entirely confined to the study of equations. In the light of the discoveries of Galois, it appears that the problem of the solution ``by radicals'' is only a particular case, fairly artificial, of the general problem of the classification of irrationals. It is this latter that, during the whole of the end of the XIXth century, is going to be attacked on several sides, and many disparate results will accumulate little by little, preparing the way for the synthesis of Steinitz.

First as far as algebraic irrationals are concerned, a fundamental principle of the classification was supplied by Galois theory, reducing the study of an algebraic equation to that of its group. In fact, it is mainly the theory of permutation groups, of which we do not have to talk here (cf.~p.~53), that, in pure Algebra, is the principal object of research at that time. Other progress in the theory of algebraic fields comes from the development, at the same time, of Number Theory and of Algebraic Geometry. This progress is concerned mainly with the way of expounding the theory, and is, for the most part, due to Dedekind {[}79{]}, who introduces the notions of field and ring, \(^{29}\) and (linked to his research on hypercomplex systems) develops systematically the linear aspect of extension theory ({[}79{]}, v. III; pp.~33 ff.). It is also he who considers the Galois group as made up of the automorphisms of the extension under consideration, and not any longer only as a group of permutations of the roots of an equation; and he proves (for the field of numbers) the fundamental theorem on the linear independence of automorphisms ({[}79{]}, v. III, p.~29) as well as the existence of normal bases for a Galois extension ({[}79{]}, v. II, p.~433). Finally, he tackles the problem of algebraic extensions of infinite degree, and notices that Galois theory can not be applied to it as it stands (an arbitrary subgroup of a Galois group not being always identical with the group of an extension with respect to a subextension); and by a bold intuition, he already begins to think of the Galois group as a topological group \(^{30}\) ---an idea that will come to fruition only with the theory of Galois extensions of infinite degree, developed by Krull in 1928 {[}187 d{]}.

Parallel to this evolution the notion of an element transcendental over a field becomes precise. The existence of transcendental numbers is proved for the first time by Liouville in 1844, by means of an explicit construction, based on the theory of diophantine approximations {[}204 c{]}; Cantor, in 1874, gives another ``non constructive'' proof using simple considerations about the power of sets {[}47{]}; finally, Hermite proves in 1873 the transcendence of e, and Lindemann in 1882 that of \(\pi\) by a method which is analogous to that of Hermite, putting thus a final seal on the ancient problem of the squaring of the circle. \(^{31}\)

As for the role of transcendental numbers in algebraic calculations, Kronecker observes in 1882 that, if x is transcendental over a field K, the field \(K(x)\) is isomorphic to the field of rational fractions \(K(X)\) ({[}186 a{]}, v. II, p.~253). He makes besides the adjunction of indeterminates to a field the keystone of his exposition of the theory of algebraic numbers ({[}186 a{]}, v. II, pp.~245-387). On the other hand, Dedekind and Weber show in the same year {[}80{]} how arithmetical methods can help to set up the theory of algebraic curves. One sees thus appearing in several directions analogies between Arithmetic and Algebraic Geometry, which will turn out to be extremely fruitful for the one and the other.

In all this research, the fields that arise consist of ``concrete'' elements in the sense of classical mathematics --- (complex) numbers or functions of complex variables. \(^{32}\) But already Kronecker, in 1882, takes full account of the fact (obscurely sensed by Gauss and Galois) that ``indeterminates'' only play the role in his theory of the basis elements of an algebra, and not that of variables in the sense of Analysis ({[}186 a{]}, v. II, pp.~339-340); and, in 1887, he develops this idea, linked to a vast programme that aims at nothing less than recasting all mathematics by rejecting all that cannot be reduced to algebraic operations on the integers. It is on this occasion that, taking up an idea of Cauchy ({[}56 a{]}, (1), v. X, pp.~312 and 351) who had defined the field C of complex numbers as the field of remainders \(\mathbf{R}[X]/(X^{2}+1)\) , Kronecker shows how the theory of algebraic numbers is completely independent of the ``fundamental theorem of algebra'' and even of the theory of real numbers, all fields of algebraic numbers (of finite degree) being isomorphic to a field of remainders \(\mathbf{Q}[X]/(f)\) (f an irreducible polynomial over \(\mathbf{Q}\) )({[}186 a{]}, v. III \(_{1}\) , pp.~211-240). As is remarked a few years later by H. Weber {[}327 a{]}, developing a first sketch of the axiomatic theory of fields, this method of Kronecker is applicable in fact to all base fields K. Weber indicates in particular that one can take as K a field \(\mathbf{Z}/(p)\) (p a prime number), thus making the calculus of congruences ``modulo p'' return to field theory; this latter had been born in the second half of the XVIIIth century, with Euler, Lagrange, Legendre and Gauss, and the analogy that it showed up with the theory of algebraic equations was not missed by observers; developing this analogy, Galois (with his research on group theory in mind) had not hesitated to introduce ``ideal roots'' of an irreducible congruence modulo p, \(^{33}\) and had indicated its principle properties ({[}123{]}, pp.~113-127). \(^{34}\) When Kronecker's method is applied to \(\mathbf{Z}/(p)\) , what is found again however (up to terminology) is the presentation that had already been given by Serret and Dedekind ({[}79{]}, v. I, p.~40) of the theory of these ``Galois imaginaries''.

To all these examples of ``abstract fields'' there come to be added, at the turn of the century, some fields of a new very different type, the fields of formal power series introduced by Veronese {[}318{]}, and above all the p-adic fields of Hensel {[}157 f{]}. It is the discovery of these latter that led Steinitz (as he says explicitly) to bring out the abstract notions common to all these theories, in a fundamental work {[}294 a{]} which can be considered as having given birth to the actual concept of Algebra. Developing systematically the consequences of the axioms of commutative fields, he introduces thus the notion of prime field, of separable (algebraic) elements, of perfect fields, defines the degree of transcendence of an extension, and proves finally the existence of algebraically closed extensions of arbitrary fields.

In recent times, the theory of Steinitz has been completed in several important areas. On the one hand, the work of Artin has brought into prominence the linear character of Galois theory {[}7 a{]}. On the other hand, the general notion of derivation (copied from the formal properties of classical differential Calculus) sensed by Dedekind ({[}79{]}, v. II, p.~412), introduced by Steinitz in the special case of a field of rational fractions ({[}294 a{]}, pp.~209-212), has been used with success in the study (essential for modern Algebraic Geometry) of transcendental extensions, and notably in the generalisation to these latter of the notion of separability {[}330 d{]}.

\chapter{DIVISIBILITY; ORDERED FIELDS.}

The elementary arithmetical operations, and especially the calculus of fractions, could not but lead to numerous empirical verifications of the divisibility of whole numbers. But neither the Babylonians (however so expert at Algebra), nor the Egyptians (in spite of their acrobatic calculations with fractions) seem to have known any general rules governing these properties, and it is to the Greeks that the initiative returns here. Their arithmetic work, of which a masterly exposition is found in Books VII and IX of Euclid {[}107{]}, does not yield pride of place to any of their most beautiful discoveries in the other branches of mathematics. The existence of the H.C.D. of two integers is proved already at the beginning of Book VII by the procedure known under the name of ``Euclid's algorithm'';¹ it serves as a foundation for all ulterior developments (properties of prime numbers, existence and calculation of the L.C.M., etc.); and the crowning of the structure is formed by the two remarkable theorems proving the existence of an infinity of prime numbers (Book IX, prop. 20) and giving a procedure for constructing perfect even numbers starting with certain prime numbers (this procedure gives in fact all the even perfect numbers, as Euler was to prove). Only the existence and uniqueness of the decomposition into prime factors are not proved in a general form; all the same Euclid proves explicitly that every integer is divisible by a prime number (Book VII, prop. 31), as well as the two following propositions (Book IX, props. 13 and 14):

``If, starting from a unit, as many numbers as desired are in a progression of constant ratio {[}i. e. geometric{]}, and that the one after the unit is prime, the largest will not be divisible by any except those that occur in the progression'' (in other words, a power \(p^n\) of a prime number can not be divided except by the powers of \(p\) with exponent \(\leq n\) ).

``If a number is the smallest that is divisible by {[}given{]} prime numbers, it will not be divisible by any other prime number with the exception of those initially {[}given as{]} dividing it'' (in other words, a product of distinct prime numbers \(p_1 \ldots p_k\) has no other prime factors than \(p_1, \ldots, p_k\) ).

It seems therefore that if Euclid does not state the general theorem it is only due to a terminology and notation adequate only for arbitrary powers of an integer. \(^{2}\)

Although a careful study makes it likely that there exists, in the text of Euclid, several successive layers, each corresponding to a stage in the development of Arithmetic, \(^{3}\) it seems that this evolution was entirely completed between the beginning of the Vth century and the middle of the IVth, and one can only admire the shrewdness and the sure footedness of the logic that are manifested there: we must wait two millennia to take part in comparable progress in Arithmetic.

It is so-called ``indeterminate'' or ``diophantine'' problems which are at the source of the ulterior developments of Number Theory. The term ``diophantine equations'', as it is used today, is not, historically, completely justified; one means in general by that algebraic equations (or systems of equations) with integral coefficients, for which one is looking only for the solutions in integers: a problem that is normally impossible if the equations are ``determined'', that is to say have only a finite number of solutions (in real or complex numbers), but that, on the contrary, often allows solutions when there are more unknowns than equations. Yet, if Diophantus seems clearly to have been the first to have considered ``indeterminate'' problems, he only exceptionally looks for integral solutions, and is content most of the time with obtaining a single solution in rational numbers {[}91 a{]}. This is a type of problem that he can solve most of the time by algebraic calculations where the arithmetic nature of the unknowns does not arise; \(^{4}\) also divisibility theory only plays a very retiring role (the word prime number is only pronounced a single time ({[}91 a{]}, Book V, problem 9, v. I, pp.~334-335), and the notion of relatively prime numbers is only brought up in connection with the theorem affirming that the quotient of two relatively prime numbers can not be a square unless each of them is a square). \(^{5}\)

The study of integral solutions of indeterminate equations does not really start until the Chinese and Hindu mathematicians of the high Middle Ages. The first of these seem to have been led to speculations of this kind by practical problems in the making of calendars (where the determination of periods common to several cycles of astronomical phenomena constituted precisely a ``diophantine'' problem of the first degree); to them is owed in any case (no doubt between the IVth and the VIIth century of our era) a rule for the solution of simultaneous linear congruences. As for the Hindus, of whom the mathematics knew its full flowering from the Vth to the XIIIth century, not only did they know how to treat methodically (by the application of Euclid's algorithm) systems of linear diophantine equations with an arbitrary number of unknowns, \(^{6}\) but they are the first to tackle and solve problems of the second degree, including certain particular cases of ``Fermat's equation'' \(Nx^{2} + 1 = y^{2}\) ({[}78{]}, vol.~II, pp.~87-307).

We do not need to pursue here the history of the theory of diophantine equations of degree \textgreater{} 1, which, following on from the work of Fermat, Euler, Lagrange and Gauss, was to end up in the XIXth century with the theory of algebraic integers (cf.~pp.~93 ff.). As we have already noted (cf.~p.~58), the study of linear systems, which seems no longer to present problems worthy of interest, is somewhat neglected during this period: in particular, general necessary conditions for solutions of an arbitrary system are not sought, nor is the description of the set of solutions. However, towards the middle of the XIXth century, Hermite is led to use, in connection with his research in Number Theory, several lemmas about linear Diophantine equations, and notably a ``reduced form'' for a linear substitution with integral coefficients ({[}159{]}, v. I, pp.~164 and 265); finally, after Heger had given in 1858 the necessary condition for a system whose rank is equal to the number of equations, H. J. Smith, in 1861, defines the invariant factors of a matrix with integral entries in a ``canonical form'' ({[}287{]}, v. I, pp.~367-409).

But in the interval the notion of an abelian group was gradually becoming more precise, following its introduction by Gauss (cf.~p.~60), and so was the importance taken by this notion in the ulterior development of Number Theory. In the particularly deep study, expounded in the Disquisitiones, of the finite abelian group of classes of quadratic forms of a given discriminant, Gauss rapidly realised that certain of these groups were not cyclic: ``in this case'', he says, ``a basis {[}that is to say a generator{]} is not sufficient, two, or a greater number, must be taken, which, by multiplication and composition, \(^{7}\) will be able to produce all the classes'' ({[}124 a{]}, v. I, pp.~374-375). It is not certain that, with these words, Gauss wished to describe the decomposition of a group as a direct product of cyclic groups; however, in the same page in the Disquisitiones, he proves that there exists an element in the group whose order is the L.C.M. of the orders of all the elements --- in other words, he obtains the existence of the largest invariant factor of the group ({[}124 a{]}, v. I, p.~373) ---; and on the other hand, the notion of direct product was known to him, for in a manuscript dating from 1801 but not published during his lifetime, he sketches a general proof of the decomposition of a finite abelian group as a direct product of \(p\)-groups \(^{8}\) ({[}124 a{]}, v. II, p.~266). In any case, in 1868, Schering, the editor of Gauss's work, inspired by these results (and notably by this manuscript that he had just found again) proves (still for the group of classes of quadratic forms) the general decomposition theorem ({[}272{]}, v. I, pp.~135-148) by a method which, taken up again two years later in abstract terms by Kronecker ({[}186 a{]}, v. I, pp.~273-282), is essentially that which is still used today. As for torsion-free abelian groups, we have already said (cf.~pp.~64 ff.) how the theory of elliptic functions and abelian integrals, developed by Gauss, Abel and Jacobi, led little by little to take note of their structure; the first and most famous example of the decomposition of an infinite group as a direct sum of monogenic groups is given by Dirichlet in 1846 in his memoir on the units of an algebraic number field ({[}92{]}, v. I, pp.~619-644). But it is only in 1879 that the link between the theory of abelian groups of finite type and the theorem of Smith is recognised and used explicitly by Frobenius and Stickelberger ({[}120{]}, §10).

At about the same time the theory of similarity for matrices (with real or complex coefficients) was being completed. The notion of eigenvalue of a linear substitution appears explicitly in the theory of systems of linear differential equations with constant coefficients, applied by Lagrange ({[}191{]}, v. I, p.~520) to the theory of small movements, by Lagrange ({[}191{]}, v. VI, pp.~655-666) and Laplace ({[}193{]}, v. VIII, pp.~325-366) to the ``secular'' inequalities of the planets. It is implicit in many other problems tackled about the middle of the XVIIIth century, such as the determination of the axes of a conic or a quadric (accomplished first by Euler ({[}108 a{]}, (1), v. IX, p.~384)), or the study (developed also by Euler ({[}108 a{]}, (2), v. III, pp.~200-201)) of the principal axes of inertia of a solid body (discovered by De Segner in 1755); we know today that it is she also who (in a much more hidden form) took part in the beginnings of the theory of partial differential equations, and in particular in the equation of vibrating strings. But (without talking of this latter case) the common ancestry between these various problems is hardly recognised before Cauchy ({[}56 a{]}, (2), v. V, p.~248 and v. IX, p.~174). Further, as the majority bring in symmetric matrices, it is the eigenvalues of these latter which are mostly studied at first; let us note here that from 1826, Cauchy proves the invariance under similarity of the eigenvalues of these matrices, and proves that they are real for a symmetric matrix of order 3 ({[}56 a{]}, (2), v. V, p.~248), a result that he generalises to arbitrary real symmetric matrices three years later ({[}56 a{]}, (2), v. IX, p.~174).⁹ The general notion of projectivity, introduced by Möbius in 1827 ({[}223{]}, v. I, p.~217), arrives quickly at the problem of classification of these transformations (in 2 and 3 dimensions first of all), which is none other than the problem of the similarity of the corresponding matrices; but for a long time this question is only treated by the ``synthetic'' methods which are in vogue in the middle of the XIXth century, and its progress (in any case fairly slow) does not seem to have had any influence on the theory of eigenvalues. It is not the same with another problem from geometry, the classification of pencils of conics or quadrics, which from the modern point of view, reduces to the study of elementary divisors of the matrix \(U + \lambda V\) where \(U\) and \(V\) are two symmetric matrices; it is exactly in this vein that Sylvester, in 1851, tackles this problem, examining with care (with the aim of finding ``canonical forms'' for the pencil under consideration) what becomes of the minors of the matrix \(U + \lambda V\) when a value is substituted for \(\lambda\) which annihilates its determinant ({[}304{]}, v. I, pp.~219-240). The purely algebraic aspect of the theory of eigenvalues progresses simultaneously; it is thus that several authors (among whom Sylvester himself) prove around 1850 that the eigenvalues of \(U^n\) are the nth powers of the eigenvalues of \(U\) , whereas in 1858 Cayley, in the memoir in which he sets up the calculus of matrices ({[}58{]}, v. II, pp.~475-496), states the ``Cayley-Hamilton theorem'' for a square matrix of arbitrary order, \(^{10}\) contenting himself with proving it by direct calculation for matrices of order 2 and 3. Finally, in 1868, Weierstrass, taking up again the method of Sylvester, obtains ``canonical forms'' for a ``pencil'' \(U + \lambda V\) , where, this time, \(U\) and \(V\) are square matrices which are not necessarily symmetric, subject only to the condition that det( \(U + \gamma V\) )

is not identically null; he deduces from this the definition of the elementary divisors of an arbitrary square matrix (with complex entries), and proves that they characterise this latter up to similarity ({[}329 a{]}, v. II, pp.~19-44); these results are besides found again partially (and apparently independently) by Jordan two years later \(^{11}\) ({[}174 a{]}, pp.~114-125). Here again, it is Frobenius who, in 1879, shows that the theorem of Weierstrass can be easily deduced from the theory of Smith, extended to polynomials ({[}119{]}, v. I, pp.~482-544, §13).

We have just alluded to the theory of divisibility of polynomials in one variable; the question of the division of polynomials was naturally to be stated from the beginning of algebra, as the inverse of the operation of multiplication (this latter being already known to Diophantus, at least for polynomials of small degree); but it can be conceived that it is hardly possible to tackle the problem in a general way before a coherent notation for the various powers of the variable was imposed. In fact, examples of the process of ``Euclidean'' division of polynomials, such as we know it, are hard to find before the middle of the XVIIth century; \(^{12}\) and S. Stévin (who essentially uses exponent notation) seems to be the first to have the idea of deducing from it the extension of ``Euclid's algorithm'' for the determination of the G.C.D. of two polynomials ({[}295{]}, v. I, pp.~54-56). Taking that into account, the notion of divisibility remained peculiar to rational integers until the middle of the XVIIIth century. It is Euler who, in 1770, opens a new chapter of Arithmetic by extending, not without temerity, the notion of divisibility to integers in a quadratic extension: seeking to determine the divisors of a number of the form \(x^{2}+cy^{2}\) ( \(x,y,c\) rational integers), he writes \(x+y\sqrt{-c}=(p+q\sqrt{-c})(r+s\sqrt{-c})\) ( \(p,q,r,s\) rational integers) and taking the norms of both members, he does not hesitate to affirm that he obtains in this way all the divisors of \(x^{2}+cy^{2}\) in the form \(p^{2}+cq^{2}\) ({[}108 a{]}, (1), v. I, p.~422). In other words, Euler argues as if the ring \(\mathbb{Z}[\sqrt{-c}]\) were principal; a little later, he uses an analogous argument to apply the method of ``infinite descent'' to the equation \(x^{3}+y^{3}=z^{3}\) (it reduces to writing that \(p^{2}+3q^{2}\) is a cube, which he does by stating \(p+q\sqrt{-3}=(r+s\sqrt{-3})^{3}\) ). But from 1773, Lagrange proves ({[}191{]}, v. III, pp.~695-795) that the divisors of numbers of the form \(x^{2}+cy^{2}\) are not always of this form, the first example of the fundamental difficulty which was to occur with much greater clarity in the studies, taken up by Gauss and his successors, on divisibility in the fields of roots of unity; \(^{13}\) it is not possible in general to extend directly to these fields the essential divisibility properties of rational integers, the existence of the G.C.D. and the uniqueness of the decomposition into prime factors. It is not the place here to describe in detail how Kummer for the field of roots of unity {[}188 b{]}, \(^{14}\) then Dedekind and Kronecker for arbitrary fields of algebraic numbers, succeed in surmounting this formidable obstacle by the creation of the theory of ideals, one of the most decisive advances of modern algebra (cf.~pp.~93 ff.). But Dedekind, always curious about the foundations of the various mathematical theories, is not content with this success; and, analysing the mechanism of the relations of divisibility, he sets up the basis of the modern theory of ordered groups, in a memoir (without stirring his contemporaries, and forgotten for 30 years) which is without doubt one of the first by date of the works on axiomatic algebra ({[}79, v. II, pp.~103-147).

From the middle of XVIIIth century, the search for a proof of the ``fundamental theorem of algebra'' is the order of the day (cf.~p.~74). We do not need to remember here the attempt by d'Alembert, who inaugurated the series of proofs using the infinitesimal calculus (cf.~p.~161). But, in 1749, Euler tackles the problem in another way ({[}108 a{]}, (1), v. VI, pp.~78-147): for every polynomial f with real coefficients, he seeks to prove the existence of a decomposition \(f = f_{1}f_{2}\) into two polynomials (not constants) with real coefficients, which would give him a proof of the ``fundamental theorem'' by induction on the degree of f.~It is sufficient even, as he remarks, to stop with the first factor of odd degree, and in consequence all the difficulties reduce to considering the case where the degree n of f is even. Euler limits himself then to the study of the case where the desired factors are both of degree n/2, and he indicates that, by an elimination calculation suitably undertaken, the unknown coefficients of \(f_{1}\) and \(f_{2}\) can be expressed rationally in terms of a root of an equation with real coefficients of which the extreme terms are of opposite signs and which as a consequence has at least one real root. But Euler's proof is only a sketch, where a number of essential points are passed over in silence; and it is only in 1772 that Lagrange succeeds in solving the difficulties raised by this proof ({[}191{]}, v. III, pp.~479-516) by means of a very long and very meticulous analysis, where he proves himself capable of a remarkable virtuosity in the use of ``Galois'' methods (so to speak) newly created by him (cf.~pp.~75 ff.).

All the same Lagrange, like Euler and all his contemporaries, does not hesitate to argue formally in a ``field of roots'' of a polynomial (that is to say, in his language, to consider ``imaginary roots'' of this polynomial); the Mathematics of his time had not supplied any justification for this type of argument. Also Gauss, deliberately hostile, from its beginnings, to the frantic formalism of the XVIIIth century, raises himself powerfully, in his dissertation, against this abuse ({[}124 a{]}, v. III, p.~3). But he would not have been himself if he had not felt that it was a case there of an argument with a defective presentation that was intrinsically correct. Also we see him, a few years later ({[}124 a{]}, v. III, p.~33; cf.~also {[}124 b{]}) taking up again a simpler variant of Euler's argument, suggested already in 1759 by de Foncenex (but that this latter was not able to bring to a satisfactory end), and deducing from it a new proof of the ``fundamental theorem'', where he avoids carefully all use of ``imaginary'' roots: these last being replaced by clever adjunctions and specialisations of indeterminates.

The role of Topology in the ``fundamental theorem'' is thus reduced to the one theorem according to which a polynomial with real coefficients cannot change sign in an interval without becoming zero (Bolzano's theorem for polynomials). This theorem is also at the bottom of all the separation criteria for real roots of a polynomial (with real coefficients), which constitute one of the favourite subjects of Algebra during the XIXth century. \(^{15}\) During this research, one could not fail to notice that it is the order structure of R, even more than its topology, that plays there the essential role; \(^{16}\) for example, Bolzano's theorem for polynomials is still true for the field of all real algebraic numbers. This movement of ideas found its goal in the abstract theory of ordered fields, created by E. Artin and O. Schreier ({[}7 b{]} and {[}8 a and b{]}); one of its more remarkable results is no doubt the discovery that the existence of an order relation on a field is linked to purely algebraic properties of the field.
